\documentclass[11pt]{article}

\usepackage{amsmath,amssymb,amsthm,bm}
\usepackage[margin=1in]{geometry}
\newtheorem{definition}{Definition}
\newtheorem{theorem}{Theorem}
\newtheorem{assumption}{Assumption}
\newtheorem{corollary}{Corollary}
\newtheorem{remark}{Remark}

\begin{document}

\section{Bounded Error of Implicit Tucker Self-Expression}
\label{app:bounded_tucker_error}

This section proves that replacing the explicit multi-view self-expression coefficient tensor with the implicit Tucker parameterization yields a bounded reconstruction error. The final error consists of four terms: ideal self-expression residual, Tucker truncation error, optimization residual, and S-PMA approximation error.

Therefore, as long as the reference self-expression tensor is approximately low Tucker-rank and the S-PMA approximation error is controlled, the implicit Tucker self-expression used by DMTSE does not cause the reconstruction error to become unbounded.

\paragraph{Notation.}
Let the multi-view traffic features be
\[
\{X^{(k)}\}_{k=1}^{K},
\qquad
X^{(k)}\in\mathbb R^{d_k\times N},
\]
where $N$ is the number of flows and $K$ is the number of views. Let
\[
\mathcal Z\in\mathbb R^{N\times N\times K}
\]
denote the reference full self-expression coefficient tensor. Its $k$-th view slice is denoted by
\[
Z^{(k)}\in\mathbb R^{N\times N}.
\]
The word ``full'' means that the tensor covers all flow pairs and all views; it does not mean that its entries are independent. In traffic data, many flows are generated by a small number of latent behavior patterns, so the full coefficient tensor can still have an approximately low Tucker-rank structure.

DMTSE uses the implicit Tucker parameterization
\[
\widehat{\mathcal Z}
=
\mathcal G\times_1 A\times_2 B\times_3 \Gamma,
\]
where $A\in\mathbb R^{N\times P}$, $
B\in\mathbb R^{N\times Q}$, $
\Gamma\in\mathbb R^{K\times R}$, and $
\mathcal G\in\mathbb R^{P\times Q\times R}$.

Let $\gamma_k$ be the $k$-th row of $\Gamma$, and define
\[
M_k
=
\mathcal G\times_3 \gamma_k
=
\sum_{r=1}^{R}\gamma_{kr}\mathcal G_{::r}
\in\mathbb R^{P\times Q}.
\]
Then the implicit coefficient matrix of view $k$ is
\[
\widehat Z^{(k)}=A M_k B^\top .
\]

Let $\phi:\mathbb R\to\mathbb R$ be the elementwise output activation in the reconstruction layer. We assume that $\phi$ is $L_\phi$-Lipschitz:
\[
|\phi(a)-\phi(b)|\le L_\phi |a-b|.
\]
For the sigmoid activation, one can take $L_\phi=1/4$; for the identity map, one can take $L_\phi=1$.

\paragraph{Spectral energy ratio.}
Let $\mathcal Z_{(n)}$ be the mode-$n$ unfolding of $\mathcal Z$, and let its singular values be sorted in non-increasing order:
\[
\sigma_1(\mathcal Z_{(n)})
\ge
\sigma_2(\mathcal Z_{(n)})
\ge
\cdots .
\]
For a target multilinear rank $(r_1,r_2,r_3)=(P,Q,R)$, define the cumulative spectral-energy ratio of the $n$-th unfolding as
\[
g_n(r_n)
=
\frac{
	\sum_{i=1}^{r_n}\sigma_i^2(\mathcal Z_{(n)})
}{
	\sum_i \sigma_i^2(\mathcal Z_{(n)})
}.
\]
Since tensor unfolding preserves the Frobenius norm,
\[
\sum_i \sigma_i^2(\mathcal Z_{(n)})
=
\|\mathcal Z\|_F^2,
\qquad n=1,2,3.
\]
Define the Tucker truncation error
\[
\varepsilon_{\rm Tuc}
=
\|\mathcal Z\|_F
\sqrt{
	\sum_{n=1}^{3}
	\left(1-g_n(r_n)\right)
},
\;
(r_1,r_2,r_3)=(P,Q,R).
\]

\begin{definition}[Optimization gap to the HOSVD witness]
	\label{def:opt_gap}
	Let $\widehat{\mathcal Z}^{\rm HO}$ be the rank-$(P,Q,R)$ truncated-HOSVD
	approximation of $\mathcal Z$, and let $\widehat{\mathcal Z}$ be the Tucker
	tensor learned by DMTSE. Define the optimization gap
	\[
	\varepsilon_{\rm opt}
	:=
	\bigl\|\widehat{\mathcal Z}-\widehat{\mathcal Z}^{\rm HO}\bigr\|_F ,
	\]
	i.e.\ the Frobenius distance between the learned factors and the
	truncated-HOSVD witness of the reference tensor. This term is small whenever
	training converges to a factorization close to the best rank-$(P,Q,R)$
	approximation; it isolates the (non-convex) optimization error from the
	representational error analyzed below.
\end{definition}

\begin{theorem}[Bounded reconstruction error of implicit Tucker self-expression]
	\label{thm:implicit_tucker_bounded_error}
	Under Definition~\ref{def:opt_gap}, define the ideal self-expression residual of view $k$ as
	\[
	\xi_k
	=
	\left\|
	X^{(k)}
	-
	\phi\!\left(X^{(k)}Z^{(k)}\right)
	\right\|_F .
	\]
	Let
	\[
	U^{(k)}=X^{(k)}A
	\]
	be the exact feature-side Tucker embedding, and let
	\[
	\widetilde U^{(k)}
	\]
	be the approximate embedding computed by S-PMA and the view-specific refiner. Define
	\[
	\eta_k
	=
	\left\|
	\widetilde U^{(k)}-U^{(k)}
	\right\|_F .
	\]
	The final DMTSE reconstruction of view $k$ is
	\[
	\widetilde X^{(k)}
	=
	\phi\!\left(
	\widetilde U^{(k)}M_kB^\top
	\right).
	\]
	
	Define the total reconstruction error
	\[
	\mathcal E_{\rm rec}
	=
	\left(
	\sum_{k=1}^{K}
	\left\|
	X^{(k)}-\widetilde X^{(k)}
	\right\|_F^2
	\right)^{1/2},
	\]
	the total ideal self-expression residual
	\[
	\varepsilon_{\rm se}
	=
	\left(
	\sum_{k=1}^{K}
	\xi_k^2
	\right)^{1/2},
	\]
	the data spectral-norm bound
	\[
	\rho_X=\max_{1\le k\le K}\|X^{(k)}\|_2,
	\]
	the Tucker view-map bound
	\[
	\rho_M=\max_{1\le k\le K}\|M_k\|_2,
	\]
	and the total S-PMA approximation error
	\[
	\varepsilon_{\rm PMA}
	=
	\left(
	\sum_{k=1}^{K}
	\eta_k^2
	\right)^{1/2}.
	\]
	Then
	\[
	\mathcal E_{\rm rec}
	\le
	\varepsilon_{\rm se}
	+
	L_\phi \rho_X
	\left(
	\varepsilon_{\rm Tuc}
	+
	\varepsilon_{\rm opt}
	\right)
	+
	L_\phi \rho_M \|B\|_2
	\varepsilon_{\rm PMA}.
	\]
	Here
	\[
	\varepsilon_{\rm Tuc}
	=
	\|\mathcal Z\|_F
	\sqrt{
		\sum_{n=1}^{3}
		\left(1-g_n(r_n)\right)
	}.
	\]
	Therefore, when the spectral energy of the three unfoldings of $\mathcal Z$ concentrates in the target ranks, the optimization gap $\varepsilon_{\rm opt}$ is small, and the S-PMA embedding error is controlled, the implicit Tucker reconstruction error is bounded.
\end{theorem}

\begin{proof}
	We restate Theorem~1 in Sec VII.A with full definitions and prove it here.
	
	\paragraph{Step 1: Tucker truncation error is controlled by tail spectral energy.}
	By the standard truncated-HOSVD approximation bound, there exists a rank-$(P,Q,R)$ Tucker tensor $\widehat{\mathcal Z}^{\rm HO}$ such that
	\[
	\left\|
	\mathcal Z-\widehat{\mathcal Z}^{\rm HO}
	\right\|_F^2
	\le
	\sum_{n=1}^{3}
	\sum_{i>r_n}
	\sigma_i^2(\mathcal Z_{(n)}).
	\]
	By the definition of $g_n(r_n)$,
	\[
	\sum_{i>r_n}
	\sigma_i^2(\mathcal Z_{(n)})
	=
	\left(1-g_n(r_n)\right)
	\|\mathcal Z\|_F^2 .
	\]
	Hence
	\[
	\left\|
	\mathcal Z-\widehat{\mathcal Z}^{\rm HO}
	\right\|_F^2
	\le
	\|\mathcal Z\|_F^2
	\sum_{n=1}^{3}
	\left(1-g_n(r_n)\right).
	\]
	Taking square roots gives
	\[
	\left\|
	\mathcal Z-\widehat{\mathcal Z}^{\rm HO}
	\right\|_F
	\le
	\varepsilon_{\rm Tuc}.
	\]
	
	\paragraph{Step 2: The optimization gap gives an explicit tensor-distance bound.}
	By Definition~\ref{def:opt_gap},
	\[
	\bigl\|\widehat{\mathcal Z}-\widehat{\mathcal Z}^{\rm HO}\bigr\|_F
	=\varepsilon_{\rm opt}.
	\]
	By the triangle inequality and Step 1,
	\[
	\left\|
	\mathcal Z-\widehat{\mathcal Z}
	\right\|_F
	\le
	\left\|
	\mathcal Z-\widehat{\mathcal Z}^{\rm HO}
	\right\|_F
	+
	\left\|
	\widehat{\mathcal Z}^{\rm HO}-\widehat{\mathcal Z}
	\right\|_F
	\le
	\varepsilon_{\rm Tuc}
	+
	\varepsilon_{\rm opt}.
	\]
	Let \(\Delta_{\mathcal Z}=\mathcal Z-\widehat{\mathcal Z}\). Then
	\[
	\|\Delta_{\mathcal Z}\|_F
	\le
	\varepsilon_{\rm Tuc}
	+
	\varepsilon_{\rm opt}.
	\]
	
	\paragraph{Step 3: Coefficient-tensor error propagates to exact Tucker reconstruction.}
	For each view $k$, define
	\[
	\Delta_k
	=
	Z^{(k)}-\widehat Z^{(k)} .
	\]
	Since the Frobenius norm of a tensor equals the square root of the sum of the squared Frobenius norms of its slices,
	\[
	\sum_{k=1}^{K}
	\|\Delta_k\|_F^2
	=
	\|\mathcal Z-\widehat{\mathcal Z}\|_F^2
	\le
	\left(
	\varepsilon_{\rm Tuc}
	+
	\varepsilon_{\rm opt}
	\right)^2 .
	\]
	
	Consider the reconstruction obtained from the learned Tucker coefficient matrix but without the S-PMA approximation:
	\[
	\widehat X_{\rm exact}^{(k)}
	=
	\phi\!\left(
	X^{(k)}\widehat Z^{(k)}
	\right).
	\]
	By the triangle inequality,
	\[
	\begin{aligned}
		\left\|
		X^{(k)}-\widehat X_{\rm exact}^{(k)}
		\right\|_F
		&=
		\left\|
		X^{(k)}
		-
		\phi\!\left(
		X^{(k)}\widehat Z^{(k)}
		\right)
		\right\|_F\\
		&\le
		\left\|
		X^{(k)}
		-
		\phi\!\left(
		X^{(k)}Z^{(k)}
		\right)
		\right\|_F\\
		&\quad+
		\left\|
		\phi\!\left(
		X^{(k)}Z^{(k)}
		\right)
		-
		\phi\!\left(
		X^{(k)}\widehat Z^{(k)}
		\right)
		\right\|_F .
	\end{aligned}
	\]
	The first term is $\xi_k$. Since $\phi$ is $L_\phi$-Lipschitz,
	\[
	\begin{aligned}
		&
		\left\|
		\phi\!\left(
		X^{(k)}Z^{(k)}
		\right)
		-
		\phi\!\left(
		X^{(k)}\widehat Z^{(k)}
		\right)
		\right\|_F\\
		&\le
		L_\phi
		\left\|
		X^{(k)}
		\left(
		Z^{(k)}-\widehat Z^{(k)}
		\right)
		\right\|_F\\
		&\le
		L_\phi
		\|X^{(k)}\|_2
		\left\|
		Z^{(k)}-\widehat Z^{(k)}
		\right\|_F\\
		&=
		L_\phi
		\|X^{(k)}\|_2
		\|\Delta_k\|_F .
	\end{aligned}
	\]
	Therefore,
	\[
	\left\|
	X^{(k)}-\widehat X_{\rm exact}^{(k)}
	\right\|_F
	\le
	\xi_k
	+
	L_\phi
	\|X^{(k)}\|_2
	\|\Delta_k\|_F .
	\]
	
	\paragraph{Step 4: S-PMA embedding error propagates to the final reconstruction.}
	By the implicit Tucker structure,
	\[
	\widehat Z^{(k)}
	=
	A M_k B^\top.
	\]
	Therefore,
	\[
	X^{(k)}\widehat Z^{(k)}
	=
	X^{(k)}A M_k B^\top
	=
	U^{(k)}M_kB^\top .
	\]
	The actual model uses the S-PMA approximation $\widetilde U^{(k)}$, so
	\[
	\widetilde X^{(k)}
	=
	\phi\!\left(
	\widetilde U^{(k)}M_kB^\top
	\right).
	\]
	Thus,
	\[
	\begin{aligned}
		&
		\left\|
		\widehat X_{\rm exact}^{(k)}
		-
		\widetilde X^{(k)}
		\right\|_F\\
		&=
		\left\|
		\phi\!\left(
		U^{(k)}M_kB^\top
		\right)
		-
		\phi\!\left(
		\widetilde U^{(k)}M_kB^\top
		\right)
		\right\|_F\\
		&\le
		L_\phi
		\left\|
		\left(
		U^{(k)}-\widetilde U^{(k)}
		\right)
		M_kB^\top
		\right\|_F\\
		&\le
		L_\phi
		\left\|
		U^{(k)}-\widetilde U^{(k)}
		\right\|_F
		\|M_k\|_2
		\|B\|_2\\
		&=
		L_\phi
		\eta_k
		\|M_k\|_2
		\|B\|_2 .
	\end{aligned}
	\]
	
	Combining Step 3 and Step 4 gives, for each view $k$,
	\[
	\left\|
	X^{(k)}-\widetilde X^{(k)}
	\right\|_F
	\le
	\xi_k
	+
	L_\phi
	\|X^{(k)}\|_2
	\|\Delta_k\|_F
	+
	L_\phi
	\eta_k
	\|M_k\|_2
	\|B\|_2 .
	\]
	Taking the squared sum over all views and applying Minkowski's inequality,
	\[
	\begin{aligned}
		\mathcal E_{\rm rec}
		&=
		\left(
		\sum_{k=1}^{K}
		\left\|
		X^{(k)}-\widetilde X^{(k)}
		\right\|_F^2
		\right)^{1/2}\\
		&\le
		\left(
		\sum_{k=1}^{K}
		\xi_k^2
		\right)^{1/2}\\
		&\quad+
		L_\phi
		\left(
		\sum_{k=1}^{K}
		\|X^{(k)}\|_2^2
		\|\Delta_k\|_F^2
		\right)^{1/2}\\
		&\quad+
		L_\phi
		\|B\|_2
		\left(
		\sum_{k=1}^{K}
		\eta_k^2
		\|M_k\|_2^2
		\right)^{1/2}.
	\end{aligned}
	\]
	By definition,
	\[
	\left(
	\sum_{k=1}^{K}
	\xi_k^2
	\right)^{1/2}
	=
	\varepsilon_{\rm se}.
	\]
	Since $\rho_X=\max_k\|X^{(k)}\|_2$,
	\[
	\left(
	\sum_{k=1}^{K}
	\|X^{(k)}\|_2^2
	\|\Delta_k\|_F^2
	\right)^{1/2}
	\le
	\rho_X
	\left(
	\sum_{k=1}^{K}
	\|\Delta_k\|_F^2
	\right)^{1/2}.
	\]
	Using Step 2,
	\[
	\left(
	\sum_{k=1}^{K}
	\|\Delta_k\|_F^2
	\right)^{1/2}
	=
	\|\mathcal Z-\widehat{\mathcal Z}\|_F
	\le
	\varepsilon_{\rm Tuc}
	+
	\varepsilon_{\rm opt}.
	\]
	Furthermore, since $\rho_M=\max_k\|M_k\|_2$,
	\[
	\left(
	\sum_{k=1}^{K}
	\eta_k^2
	\|M_k\|_2^2
	\right)^{1/2}
	\le
	\rho_M
	\left(
	\sum_{k=1}^{K}
	\eta_k^2
	\right)^{1/2}
	=
	\rho_M
	\varepsilon_{\rm PMA}.
	\]
	Therefore,
	\[
	\mathcal E_{\rm rec}
	\le
	\varepsilon_{\rm se}
	+
	L_\phi \rho_X
	\left(
	\varepsilon_{\rm Tuc}
	+
	\varepsilon_{\rm opt}
	\right)
	+
	L_\phi \rho_M \|B\|_2
	\varepsilon_{\rm PMA}.
	\]
	Substituting the expression of $\varepsilon_{\rm Tuc}$ completes the proof.
\end{proof}

\begin{corollary}[Explicit bound induced by per-position S-PMA error]
	\label{cor:spma_to_recon}
	Suppose S-PMA satisfies, for every view $k$ and every feature position $m$,
	\[
	\left\|
	\widetilde u_m^{(k)}-u_m^{(k)}
	\right\|_2
	\le
	C_{\rm spma}^{(k)} .
	\]
	Then
	\[
	\eta_k
	=
	\left\|
	\widetilde U^{(k)}-U^{(k)}
	\right\|_F
	\le
	\sqrt{d_k}\,C_{\rm spma}^{(k)}.
	\]
	Therefore,
	\[
	\varepsilon_{\rm PMA}
	\le
	\left(
	\sum_{k=1}^{K}
	d_k
	\left(C_{\rm spma}^{(k)}\right)^2
	\right)^{1/2}.
	\]
	If all views share a common bound $C_{\rm spma}$, then
	\[
	\varepsilon_{\rm PMA}
	\le
	C_{\rm spma}
	\sqrt{
		\sum_{k=1}^{K}d_k
	}.
	\]
	Substituting it into Theorem~\ref{thm:implicit_tucker_bounded_error} gives
	\[
	\mathcal E_{\rm rec}
	\le
	\varepsilon_{\rm se}
	+
	L_\phi \rho_X
	\left(
	\varepsilon_{\rm Tuc}
	+
	\varepsilon_{\rm opt}
	\right)
	+
	L_\phi \rho_M\|B\|_2
	C_{\rm spma}
	\sqrt{
		\sum_{k=1}^{K}d_k
	}.
	\]
\end{corollary}

\begin{proof}
	By the definition of the Frobenius norm,
	\[
	\eta_k^2
	=
	\left\|
	\widetilde U^{(k)}-U^{(k)}
	\right\|_F^2
	=
	\sum_{m=1}^{d_k}
	\left\|
	\widetilde u_m^{(k)}-u_m^{(k)}
	\right\|_2^2 .
	\]
	If
	\[
	\left\|
	\widetilde u_m^{(k)}-u_m^{(k)}
	\right\|_2
	\le
	C_{\rm spma}^{(k)},
	\]
	then
	\[
	\eta_k^2
	\le
	d_k
	\left(C_{\rm spma}^{(k)}\right)^2,
	\]
	which gives
	\[
	\eta_k
	\le
	\sqrt{d_k}\,C_{\rm spma}^{(k)}.
	\]
	Summing over all views yields
	\[
	\varepsilon_{\rm PMA}
	=
	\left(
	\sum_{k=1}^{K}
	\eta_k^2
	\right)^{1/2}
	\le
	\left(
	\sum_{k=1}^{K}
	d_k
	\left(C_{\rm spma}^{(k)}\right)^2
	\right)^{1/2}.
	\]
	The shared-bound case follows immediately by setting $C_{\rm spma}^{(k)}=C_{\rm spma}$. Substitution into Theorem~\ref{thm:implicit_tucker_bounded_error} completes the proof.
\end{proof}

\begin{remark}[Interpretation]
	Theorem~\ref{thm:implicit_tucker_bounded_error} shows that the final reconstruction error of DMTSE is controlled by four terms:
	\[
	\varepsilon_{\rm se},
	\qquad
	\varepsilon_{\rm Tuc},
	\qquad
	\varepsilon_{\rm opt},
	\qquad
	\varepsilon_{\rm PMA}.
	\]
	The term $\varepsilon_{\rm se}$ measures how well the original data can be explained by an ideal full self-expression model. The term $\varepsilon_{\rm Tuc}$ is determined by the tail singular-value energy of the three unfoldings of the full coefficient tensor, and is directly supported by the mode-wise spectral-energy measurements. The term $\varepsilon_{\rm opt}$ makes the optimization gap to the truncated-HOSVD witness explicit. The term $\varepsilon_{\rm PMA}$ measures the approximation error introduced when S-PMA replaces the full feature-side aggregation $U^{(k)}=X^{(k)}A$. Therefore, if the traffic self-expression tensor has a clear low Tucker-rank structure and S-PMA gives a controlled approximation, the implicit Tucker generator can avoid materializing the $N\times N\times K$ tensor while maintaining a bounded reconstruction error.
\end{remark}

\end{document}

